\title{\Large\textbf{Building a Real-Time Data Lakehouse Architecture for Intelligent Transportation Systems on a Kubernetes Cluster}}
\author{
    Ngô Tiến Sỹ (23521367) \quad Nguyễn Văn Nam (23520982) \\[0.5em]
    \fontsize{10pt}{11pt}\selectfont \textit{University of Information Technology -- VNU HCMC} \\
    \fontsize{10pt}{11pt}\selectfont \textit{Distributed Databases and Big Data -- IS211.Q22 \& IS405.Q23} \\
    \fontsize{10pt}{11pt}\selectfont \textit{Supervisor: M.Sc. Nguyễn Hồ Duy Trí}
}
\date{Ho Chi Minh City, 2026}

\maketitle

\begin{abstract}
This study builds and evaluates \textbf{Continux}, a real-time Data Lakehouse architecture for intelligent transportation data on a three-server K3s cluster. The system converts NYC TLC Yellow Taxi data from Parquet format to JSONL format, then replays the records as an event stream with Vector, sends the events into Redpanda as an intermediate layer for stream ingestion and distribution, processes them with streaming SQL queries in RisingWave, and writes the results into Apache Iceberg tables stored on MinIO; at the same time, deployment is managed by Argo CD and system behavior is observed through VictoriaMetrics/Grafana. The artifact (the concrete experimental system and its components) is evaluated through one controlled replay and a materialized-view cutover scenario between the old and new branches: the public MV (the materialized view exposed for external queries) is switched to the green logic; the query loop records no errors, RisingWave does not restart, and the system records a \texttt{0.145226 s} cutover duration. The result shows that public query logic can be updated without observed interruption in the measured window, while the pipeline also produces observable Iceberg output on MinIO.
\end{abstract}

\noindent\textbf{Keywords:} Data Lakehouse, stream processing, Materialized View, Blue/Green cutover, RisingWave, Apache Iceberg, GitOps, K3s, ITS.
